% Copyright 2026 Chris Payne
%
% Permission is hereby granted, free of charge, to any person
% obtaining a copy of this software and associated documentation files
% (the "Software"), to deal in the Software without restriction,
% including without limitation the rights to use, copy, modify, merge,
% publish, distribute, sublicense, and/or sell copies of the Software,
% and to permit persons to whom the Software is furnished to do so,
% subject to the following conditions:
%
% The above copyright notice and this permission notice shall be
% included in all copies or substantial portions of the Software.
%
% THE SOFTWARE IS PROVIDED "AS IS", WITHOUT WARRANTY OF ANY KIND,
% EXPRESS OR IMPLIED, INCLUDING BUT NOT LIMITED TO THE WARRANTIES OF
% MERCHANTABILITY, FITNESS FOR A PARTICULAR PURPOSE AND
% NONINFRINGEMENT. IN NO EVENT SHALL THE AUTHORS OR COPYRIGHT HOLDERS
% BE LIABLE FOR ANY CLAIM, DAMAGES OR OTHER LIABILITY, WHETHER IN AN
% ACTION OF CONTRACT, TORT OR OTHERWISE, ARISING FROM, OUT OF OR IN
% CONNECTION WITH THE SOFTWARE OR THE USE OR OTHER DEALINGS IN THE
% SOFTWARE.

\chapter {Distribution}\label{chap:dist}

\section {Introduction}

The distribution library\index{distribution} lets processes in
different Swish instances, called \emph{nodes}\index{node}, send
messages to each other with the same \code{send} procedure used for
local processes. A process on another node is represented locally by a
\emph{remote process identifier}\index{remote process identifier}, or
remote pid. Remote pids satisfy \code{process?}, and \code{send},
\code{link}, \code{unlink}, \code{monitor}, \code{demonitor},
\code{demonitor\&flush}, and \code{kill} accept them. As a result,
higher-level behaviors such as \code{gen-server:call} and supervisors
work across nodes without change.

A node starts distribution by calling \code{dist:start\&link} with a
node name and a shared secret, called a \emph{cookie}\index{cookie}.
Nodes that use the same cookie can connect to each other with
\code{dist:connect}. Each pair of nodes shares one connection, over TCP
or WebSocket. Once two nodes are connected, a message that contains
processes or monitors can be sent from one node to the other, and the
receiver gets the same identities back.

The following example connects two nodes on the same machine. On node
\code{alpha}, a process registered as \code{echo} replies to each
message:\antipar

\begin{alltt}
(dist:start\&link 'alpha@localhost "secret" 9001 "127.0.0.1")
(register 'echo
  (spawn
   (lambda ()
     (let lp ()
       (receive
        [#(echo ,from ,msg) (send from `#(echoed ,msg))])
       (lp)))))
\end{alltt}

On node \code{beta}, a process connects to \code{alpha}, sends a
message to the registered name \code{echo} on node
\code{alpha@localhost}, and waits for the reply:\antipar

\begin{alltt}
(dist:start\&link 'beta@localhost "secret" 9002 "127.0.0.1")
(dist:connect "127.0.0.1:9001")         \(\Rightarrow\) #(ok alpha@localhost)
(send '#(echo alpha@localhost) `#(echo ,self "hello"))
(receive [#(echoed ,msg) msg])          \(\Rightarrow\) "hello"
\end{alltt}

The library also provides remote spawn
(Section~\ref{sec:dist-remote-spawn}) and a name registry shared by all
connected nodes (Section~\ref{sec:dist-global-names}).

The design follows Erlang's distribution, with these
differences:\antipar

\begin{itemize}
\item Nodes connect only when \code{dist:connect} is called. Sending to
  an unconnected node does not connect to it, and connected nodes do
  not form a mesh.
\item Sending a message that cannot be serialized, such as one that
  contains a procedure, raises an exception in the sender.
\item There is no heartbeat. A node notices a lost connection when the
  underlying socket closes.
\end{itemize}

\section {Theory of Operation}

\subsection {Nodes}

A node name\index{node!name} is a symbol of the form
\code{\var{name}@\var{host}}. When \code{dist:start\&link} is given a
name without \code{@}, it appends \code{@} and the host name from
\code{osi\_get\_hostname}. A node also has a
\emph{creation}\index{node!creation}, the value of \code{(erlang:now)}
when distribution started. A process identity on the wire includes the
node name, creation, and process id. Remote pids from an earlier run of
a node with the same name therefore never resolve to processes in the
new run.

Without distribution, \code{(dist:node)} returns
\code{nonode@nohost}.

\subsection {Remote Process Identifiers}

A remote pid is a record with a node name, a creation, and a process id
on that node. The distribution library interns remote pids, so a remote
process has a single remote pid on each node. Remote pids can therefore
be compared with \code{eq?} and matched by \code{match} patterns such as
\code{,@pid}. Remote monitors are interned the same way, so the
\code{,@m} patterns used by \code{gen-server:call} work across nodes.

The process procedures of Chapter~\ref{chap:erlang} treat remote pids as
follows:\antipar

\begin{itemize}
\item \code{process?} returns \code{\#t}.
\item \code{process-id} returns the process's id on its own node, so it
  can equal the id of an unrelated local process.
\item \code{process-name} returns \code{\#f}, even if the process is
  registered on its own node.
\item \code{global-process-id} returns
  \code{"\var{node}:\var{id}"}.
\item A remote pid prints as \code{\#<process \var{node}:\var{id}>}.
\end{itemize}

A local process that is sent to another node, directly or inside a
message, is \emph{exported}\index{process!exported}. The distribution
library holds a strong reference to an exported process until the
process dies, so that remote nodes can still reach it. As a result,
Swish does not garbage collect a blocked process once its pid has been
sent to another node. After the process dies, the library holds it
weakly, so that its identity resolves to the same local object for as
long as local references to it remain.

\subsection {Operations on Remote Processes}

Operations on a remote pid are forwarded to its node over the
connection to that node. Operations between one pair of processes
arrive in the order they were sent.\antipar

\begin{itemize}
\item \code{(send \var{rpid} \var{message})} serializes \var{message}
  in the sender. If serialization fails, the exception is raised in the
  sender. If the node is not connected, the message is discarded. The
  receiving node discards a message for a process that has exited.

\item \code{(link \var{rpid})} records the link locally and asks the
  remote node to record it too. If the remote process has exited, the
  caller receives an exit signal with reason \code{noproc}. If the node
  is not connected, the reason is \code{noconnection}. As with a local
  link, the caller receives an \code{EXIT} message if it traps exits,
  and otherwise it is killed.

\item \code{(monitor \var{rpid})} returns a monitor. If the remote
  process has exited, the caller receives a \code{DOWN} message with
  reason \code{noproc}. If the node is not connected, the reason is
  \code{noconnection}.

\item \code{(kill \var{rpid} \var{reason})} sends an exit signal that
  the remote node handles as a local \code{kill} would: reason
  \code{kill} kills the process with reason \code{killed}, a process
  that traps exits receives an \code{EXIT} message, and reason
  \code{normal} is otherwise ignored.

\item \code{unlink}, \code{demonitor}, and \code{demonitor\&flush}
  remove the link or monitor on both nodes.
\end{itemize}

When a local process exits, the library sends an exit signal to each
linked remote process and a \code{DOWN} message to each remote process
that monitors it. Exit reasons are sent without their fault-condition
wrapper, so the stack is not available on the receiving node. A reason
that cannot be serialized, such as a Chez Scheme condition, is replaced
by \code{\#(unserializable-reason \var{string})}, where \var{string} is
the printed reason.

When a connection closes, every link to a process on that node fires
with reason \code{noconnection}, and every monitor of a process on that
node delivers a \code{DOWN} message with reason
\code{noconnection}.

When distribution is not running, a remote pid acts like a process on a
node that is not connected: \code{send} discards the message,
\code{link} delivers an exit signal with reason \code{noconnection}, and
\code{monitor} delivers a \code{DOWN} message with reason
\code{noconnection}.

\subsection {Destinations}

In addition to processes and registered names, \code{send} accepts the
following destinations when distribution is running:\antipar

\begin{itemize}
\item \code{\#(\var{name} \var{node})}: the process registered as
  \var{name} on \var{node}. If \var{node} is the local node, this is
  the same as sending to \var{name}. If \var{node} is not connected,
  the message is discarded. If no process is registered as \var{name}
  on \var{node}, the receiving node discards the message.

\item \code{\#(global \var{name})}: the process registered as
  \var{name} with \code{global:register}. If there is no such process,
  exception \code{\#(bad-arg send \var{destination})} is raised. This
  form takes precedence over \code{\#(\var{name} \var{node})}, so a
  process registered locally as \code{global} cannot be reached as
  \code{\#(global \var{node})}.
\end{itemize}

\subsection {Connections}

A connection carries \emph{frames}\index{distribution!frame}. A frame
is a 4-byte big-endian length followed by two objects encoded with
\code{fasl-write}: a vector of tokens and a body. When writing the
body, the library uses \code{fasl-write}'s external-object predicate
to replace each process and monitor with an index into the token
vector. A process token is \code{\#(p \var{node} \var{creation}
\var{id})}. A monitor token is \code{\#(m \var{key} \var{origin}
\var{target})}, where \var{key} is \code{\#(\var{node}
\var{creation} \var{id})} for the node that created the monitor and
\var{origin} and \var{target} are process tokens. The receiver
converts the tokens to local processes or interned remote pids and
monitors, and it passes the result to \code{fasl-read} as the
externals vector. A frame is limited to 64~MiB.

The TCP carrier writes frames to a socket with \code{TCP\_NODELAY}
set. The WebSocket carrier sends each frame, including its length, as
one binary WebSocket message.

\paragraph* {handshake}\index{distribution!handshake} After a
connection opens, both sides send \code{\#(hello \var{version}
\var{node} \var{creation} \var{challenge})}, where \var{challenge} is a
random 16-byte bytevector. Each side then sends \code{\#(proof
\var{proof})}, where \var{proof} is the HMAC-SHA1 of the peer's
challenge followed by the UTF-8 node name of the sender, keyed with the
cookie. A connection is closed if the peer's proof does not match
(reason \code{bad-cookie}), if the peer uses the local node name
(reason \code{same-node}), if the handshake does not finish within 10
seconds (reason \code{handshake-timeout}), or if the node is already
connected. In the last case, the existing connection is kept.

Each connection runs in a process that writes outgoing frames, and a
linked reader process decodes and handles incoming frames. Both
processes exit when the carrier fails, and the connection process
closes the carrier when it exits.

\subsection {Supervision}

\code{dist:start\&link} starts a one-for-all supervisor registered as
\code{dist-sup} with the following children:\antipar

\begin{itemize}
\item \code{dist-conn-sup}: a watcher for the connection processes
\item \code{net-kernel}: the gen-server that listens for and starts
  connections
\item \code{global-name-server}: the gen-server that holds this node's
  copy of the global name table
\item \code{dist-http}: an HTTP server that accepts WebSocket
  connections, present only when a WebSocket port is given
\end{itemize}

The supervisor allows no restarts, so distribution stops if any child
fails.

\genserver{net-kernel}{init} The \code{init} procedure traps exits,
listens for TCP connections on the given address and port, records the
node name and a new creation, and installs the hooks through which
\code{send}, \code{link}, \code{monitor}, and related procedures reach
the distribution library. Its state holds the listener, the cookie,
the open connections, and the callers waiting in
\code{dist:connect}.

\genserver{net-kernel}{terminate} The \code{terminate} procedure closes
the listener and removes the hooks. Every link and monitor that refers
to a remote process then fires with reason \code{noconnection}.

\genserver{net-kernel}{handle-call} The \code{handle-call} procedure
handles the following messages:\antipar

\begin{itemize}
\item \code{\#(connect \var{address})}: Start a connection to
  \var{address} under \code{dist-conn-sup}, and reply when it either
  finishes its handshake or exits.
\item \code{\#(up \var{node} \var{creation} \var{pid})}: Sent by a
  connection after its handshake. If \var{node} is already connected,
  reply \code{duplicate}. Otherwise, record \var{pid} as the connection
  for \var{node}, notify \code{global-name-server}, and reply
  \code{ok}.
\item \code{\#(accept-ws \var{handler})}: Start a connection for a
  WebSocket request being handled by \var{handler}.
\item \code{listen-port}: Reply with the TCP listener's port number.
\item \code{\#(disconnect \var{node})}: Kill the connection for
  \var{node} with reason \code{shutdown}, if any.
\end{itemize}

\genserver{net-kernel}{handle-info} The \code{handle-info} procedure
starts a connection for each accepted TCP connection. It stops when the
listener fails. When a connection exits, it removes the node and
delivers \code{noconnection} for that node's links and monitors. When
an exported process exits, it switches that process to a weak
reference.

\subsection {Remote Spawn}\label{sec:dist-remote-spawn}

\code{dist:spawn} and \code{dist:spawn\&link} start a process on
another node. No code crosses the wire. The caller names a procedure
with \code{\#(\var{module} \var{proc} \var{args})}, and the remote node
looks it up with \code{(top-level-value \var{proc} (environment
\var{module}))}, or in \code{(interaction-environment)} when
\var{module} is \code{\#f}. The procedure must therefore be defined on
the remote node. The remote node applies it to \var{args} in a new
process. For \code{dist:spawn\&link}, the link is set up on both nodes
before the caller receives the pid, so an early exit is not missed.

\subsection {Global Names}\label{sec:dist-global-names}

Global names\index{global names} are names registered on all connected
nodes. Each node's \code{global-name-server} holds a copy of the name
table. To register a name, the caller locks it on every
\code{global-name-server}, in order of node name, and then registers it
on each one, which also releases the lock. If a lock is held by another
caller, the caller releases the locks it holds, waits a random 1--20
milliseconds, and tries again. Each server monitors registered
processes and lock holders, so a name disappears when its process exits
or its node disconnects.

When two nodes connect, their servers exchange tables. If both have a
process registered under the same name, for example after a network
partition heals, each server keeps the process that sorts first by
node name and then by process id. It kills the other process with
reason \code{\#(global-name-conflict \var{name})}. Both nodes choose the
same winner.

Global names reach only nodes that are directly connected. A name
registered between a connect and the table exchange is not locked
against the new node, and a conflict of this kind is resolved in the
same way.

\subsection {WebSocket Transport}

A node accepts WebSocket connections at path \code{/dist} when
\code{dist:start\&link} is given a WebSocket port, or when an
application adds \code{(dist:url-handler)} to its own HTTP server (see
Chapter~\ref{chap:http}). An address of the form
\code{ws://\var{host}:\var{port}/\var{path}} makes \code{dist:connect}
use WebSocket. The handshake and everything above it are the same as
for TCP.

\subsection {Security}

The cookie handshake keeps nodes with different cookies from
connecting, but it is not encryption: frames are sent in the clear. A
connected node is fully trusted. It can send any object that
\code{fasl-read} can build, kill local processes, and call any
procedure that \code{dist:spawn} can name. Use distribution only on
trusted networks, or carry it over a secure tunnel.

\subsection {Limitations}

\begin{itemize}
\item If two nodes connect to each other at the same moment, both
  connections can be closed as duplicates. Calling \code{dist:connect}
  again establishes a single connection.
\item When a connection closes, the library scans every process to
  deliver \code{noconnection}.
\item A link exit signal for a link that no longer exists is ignored.
  There is no unlink acknowledgment.
\item The WebSocket server started by \code{dist:start\&link} listens on
  all interfaces, regardless of the address argument.
\end{itemize}

\section {Programming Interface}

\defineentry{dist:start\&link}
\begin{procedure}
  \code{(dist:start\&link \var{name} \var{cookie} \opt{\var{port}
      \opt{\var{address} \opt{\var{ws-port}}}})}
\end{procedure}
\returns{}
\code{\#(ok \var{pid})} $|$
\code{\#(error \var{reason})}

The \code{dist:start\&link} procedure starts distribution under a
supervisor registered as \code{dist-sup} and links it to the calling
process. \var{name} is a symbol. If it does not contain \code{@},
\code{@} and the host name are appended. \var{cookie} is a string,
which is converted to UTF-8, or a bytevector. The node listens for TCP
connections on \var{port} (default 0, for any available port) at
\var{address} (default \code{"0.0.0.0"}). If \var{ws-port} is not
\code{\#f} (the default), the node also accepts WebSocket connections
at \code{/dist} on an HTTP server listening on \var{ws-port}.

Because the supervisor is linked to the caller, a process that exits
while linked to it stops distribution. Start distribution from a
long-lived process or from an application's supervision tree.

\defineentry{dist:stop}
\begin{procedure}
  \code{(dist:stop)}
\end{procedure}
\returns{} \code{ok}

The \code{dist:stop} procedure kills \code{dist-sup} with reason
\code{shutdown}, if it is running, and waits for it to exit. All
connections close, so links and monitors to remote processes fire with
reason \code{noconnection}.

\defineentry{dist:connect}
\begin{procedure}
  \code{(dist:connect \var{address} \opt{\var{timeout}})}
\end{procedure}
\returns{}
\code{\#(ok \var{node})} $|$
\code{\#(error \var{reason})}

The \code{dist:connect} procedure connects to the node at \var{address}
and returns its name. \var{address} is a string of the form
\code{"\var{host}:\var{port}"}, \code{"[\var{ipv6-host}]:\var{port}"},
or \code{"ws://\var{host}:\var{port}\opt{/\var{path}}"}, where
\var{path} defaults to \code{/dist}. If the address cannot be parsed,
exception \code{\#(bad-address \var{address})} is raised. If the
connection or handshake fails, \var{reason} is the connection's exit
reason, such as \code{bad-cookie}. \var{timeout} is passed to
\code{gen-server:call} and defaults to 30000 milliseconds. If \var{node}
is already connected, the existing connection is kept and
\code{\#(ok \var{node})} is returned.

\defineentry{dist:disconnect}
\begin{procedure}
  \code{(dist:disconnect \var{node})}
\end{procedure}
\returns{} a boolean

The \code{dist:disconnect} procedure closes the connection to
\var{node} and returns \code{\#t}. If \var{node} is not connected, it
returns \code{\#f}. Links and monitors to processes on \var{node} fire
with reason \code{noconnection}.

\defineentry{dist:node}
\begin{procedure}
  \code{(dist:node \opt{\var{process}})}
\end{procedure}
\returns{} a symbol

The \code{dist:node} procedure returns the name of the node on which
\var{process} runs. Without an argument, it returns the local node
name, or \code{nonode@nohost} when distribution is not running. If
\var{process} is not a process, exception \code{\#(bad-arg dist:node
\var{process})} is raised.

\defineentry{dist:nodes}
\begin{procedure}
  \code{(dist:nodes)}
\end{procedure}
\returns{} a list of symbols

The \code{dist:nodes} procedure returns the names of the connected
nodes, sorted alphabetically.

\defineentry{dist:listen-port}
\begin{procedure}
  \code{(dist:listen-port)}
\end{procedure}
\returns{} a port number

The \code{dist:listen-port} procedure returns the TCP port on which the
node accepts connections. It is useful when \code{dist:start\&link}
was called with port 0.

\defineentry{dist:websocket-port}
\begin{procedure}
  \code{(dist:websocket-port)}
\end{procedure}
\returns{} a port number

The \code{dist:websocket-port} procedure returns the port of the HTTP
server started by \code{dist:start\&link} for WebSocket connections.

\defineentry{dist:url-handler}
\begin{procedure}
  \code{(dist:url-handler)}
\end{procedure}
\returns{} a URL handler

The \code{dist:url-handler} procedure returns an
\code{http:url-handler} that accepts distribution connections over
WebSocket at path \code{/dist}. Use it to accept connections on an
application's existing HTTP server. The handler declines requests for
other paths, and all requests when distribution is not running.

\defineentry{dist:spawn}
\begin{procedure}
  \code{(dist:spawn \var{node} \#(\var{module} \var{proc} \var{args}))}
\end{procedure}
\returns{} a process

\defineentry{dist:spawn\&link}
\begin{procedure}
  \code{(dist:spawn\&link \var{node} \#(\var{module} \var{proc} \var{args}))}
\end{procedure}
\returns{} a process

The \code{dist:spawn} procedure creates a process on \var{node} that
applies the procedure named \var{proc} to the list \var{args}, and
returns its pid. \var{module} is a library name, such as \code{(my
lib)}, or \code{\#f} for the interaction environment, and \var{proc} is
a symbol. If \var{node} is the local node, the process is created
locally. The \code{dist:spawn\&link} procedure also links the new
process to the caller, before the pid is returned.

The following exceptions can be raised:\antipar

\begin{itemize}
\item \code{\#(noconnection \var{node})}: \var{node} is not connected,
  or the connection closed before the reply arrived.
\item \code{\#(bad-module \var{module})}: \var{module} cannot be found
  on \var{node}.
\item \code{\#(undefined-procedure \var{module} \var{proc})}:
  \var{proc} is not bound in \var{module}.
\item \code{\#(not-a-procedure \var{module} \var{proc})}: \var{proc}
  is bound to a value that is not a procedure.
\end{itemize}

\defineentry{global:register}
\begin{procedure}
  \code{(global:register \var{name} \var{process})}
\end{procedure}
\returns{} \code{\#t}

The \code{global:register} procedure registers \var{process} as
\var{name} on every connected node. \var{name} is a symbol, and
\var{process} may be local or remote. If \var{name} is already
registered, exception \code{\#(name-already-registered \var{pid})} is
raised. If the name stays locked by other callers after about 100
attempts, exception \code{\#(global-lock-timeout \var{name})} is
raised.

\defineentry{global:unregister}
\begin{procedure}
  \code{(global:unregister \var{name})}
\end{procedure}
\returns{} \code{\#t}

The \code{global:unregister} procedure removes \var{name} from every
connected node.

\defineentry{global:whereis}
\begin{procedure}
  \code{(global:whereis \var{name})}
\end{procedure}
\returns{} a process or \code{\#f}

The \code{global:whereis} procedure returns the process registered
globally as \var{name}, or \code{\#f} if there is none.
